\section{Results}
\label{sec:results}

All headline results use frozen predictors and paired replayed request traces. Unless a table or caption explicitly labels a quantity as a predictor diagnostic, reported values are realized simulator outcomes. The strengthened evidence program separates approximation quality, filtering behavior, uptake-regime sensitivity, and external directional checks so that no claim is stronger than its evidence tier (Table~\ref{tab:evidence_hierarchy_main}).

\subsection{Exact-vs-Greedy Assortment Diagnostics}
\label{sec:exact_greedy}

\textit{[Mechanism diagnostic.]} This section provides mechanism diagnostics for the menu-construction approximation.

\ArtifactTable{../artifacts/tables/phase29_exact_greedy_gap_summary.tex}{Exact-vs-greedy assortment diagnostics on the RC benchmark.}

Table~\ref{tab:phase29_exact_greedy_gap} reports the main approximation-quality diagnostic. Exact enumeration is interpreted as the benchmark for the implemented menu-value surrogate when the candidate set is small enough to enumerate. Greedy forward selection is the scalable approximation. The logged relative optimality gap, menu-overlap rate, and build-time columns therefore answer a narrow but important question: how much surrogate value is lost when the system uses the operationally cheaper greedy construction rather than exact subset enumeration? This table is the primary evidence for the exact-small/greedy-large algorithmic claim; it should not be read as a proof of optimality for the full dynamic DRT routing problem.

\subsection{Robust ETA Filtering}
\label{sec:robust_filtering}

\textit{[Mechanism diagnostic.]}

\ArtifactTable{../artifacts/tables/phase30_robust_filtering_summary.tex}{Robust ETA-filtering comparison on the RC benchmark.}

Table~\ref{tab:phase30_robust_filtering} compares hard ETA filtering, calibrated filtering, interval-overlap filtering, and no ETA filtering. The key diagnostic is not only net profit, but also false-negative pruning, home-only menu share, and displayed menu size. A robust filter is useful only if it preserves operational credibility without removing most feasible meeting-point bundles before assortment optimization acts. The no-filter variant is included as an upper-bound diagnostic on candidate availability, not as a universal recommendation.

\subsection{Uptake-Regime Menu Value}
\label{sec:uptake_regime}

\textit{[Behavioral stress test.]}

\ArtifactTable{../artifacts/tables/phase31_uptake_menu_value_summary.tex}{Menu-value comparison across low, medium, and high uptake regimes.}

Table~\ref{tab:phase31_uptake_menu_value} reports the low-, medium-, and high-uptake RC regimes. This table is the main guard against over-interpreting near-degenerate outside-option settings. When acceptance and non-home uptake are close to zero, profit differences mainly diagnose mechanism behavior; when uptake becomes behaviorally live, the same menu and pricing rules can be interpreted through operator profit, passenger surplus, non-home acceptance, and pricing-boundary diagnostics. The intended claim is therefore conditional: service-menu optimization matters most when passengers have enough willingness to consider meeting-point bundles and the outside option is not overwhelmingly attractive.

\subsection{MNL Parameter Sensitivity}
\label{sec:mnl_sensitivity}

\textit{[Behavioral stress test.]}

\ArtifactTable{../artifacts/tables/mnl_sensitivity_summary.tex}{MNL-regime sensitivity summary across low, medium, and higher-uptake RC regimes.}

Table~\ref{tab:mnl_sensitivity_summary} reports the MNL parameter sensitivity results across low-, medium-, and higher-uptake regimes. This evidence addresses the concern that the menu-optimization results are artifacts of a single MNL calibration. The base utility, home-pickup bonus, price sensitivity, walking scale, and outside-option strength are all embedded in the three regime configurations. When acceptance and non-home uptake remain near zero (low regime), profit differences are mechanism diagnostics; when uptake increases (medium and higher regimes), the same menu and pricing rules produce richer behavioral variation.

The sensitivity evidence shows that the direction and rough magnitude of the menu-optimization gap are stable across the tested MNL parameterizations. This is a stress-test result conditional on the implemented parameterization, not a claim that results generalize to all possible demand curves. An explicit outside-option utility scan would require code changes to the choice model and is deferred to future work.

\subsection{RC Outside-Option Benchmark}

\textit{[Mechanism diagnostic.]}

\ArtifactTable{../artifacts/tables/rc_main_optout_summary.tex}{Main RC policy comparison under the outside-option evaluation setup. Gap and win rate are computed against exhaustive full display across six split-level evaluation pairs.}
\ArtifactFigure{../artifacts/figures/rc_main_optout_gap_vs_full.png}{0.9\linewidth}{Policy net-profit gap against full display on the RC outside-option study.}{fig:main_policy_gap}

The original RC outside-option benchmark remains useful as a mechanism check. It shows whether policies change menu exposure, home-only collapse, and candidate-set pruning under a demanding outside-option calibration. It is not the main operator-impact benchmark when full-display acceptance is behaviorally degenerate. Split-level uncertainty for RC, Austin, and Seattle is reported in Table~\ref{tab:policy_gap_uncertainty_summary}. The RC benchmark provides six split-level evaluation pairs. Reported gaps are split-level paired means with the observed range; the percentage effect is computed relative to the full-display net profit. When acceptance rates differ between policies (e.g., strict filtering reduces non-home acceptance), the profit gap partially reflects changed request composition rather than purely operational efficiency.

\subsection{Cross-Instance Evaluation}
\label{sec:cross_instance}

\textit{[Descriptive external check.]} Austin and Seattle each provide only two split-level evaluation pairs; the city results are reported as descriptive two-pair checks, not as stable generalization estimates.

\ArtifactTable{../artifacts/tables/austin_main_summary.tex}{Austin (Amazon Last Mile) policy comparison under the outside-option evaluation setup. Results averaged over two split-level evaluation pairs ($0{\to}1$, $1{\to}0$).}

\ArtifactTable{../artifacts/tables/seattle_main_summary.tex}{Seattle (Amazon Last Mile) policy comparison under the outside-option evaluation setup. Results averaged over two split-level evaluation pairs ($0{\to}1$, $1{\to}0$).}

\ArtifactTable{../artifacts/tables/city_impact_summary.tex}{City-level management summary for the no-filter heuristic against full display and the strict-filter heuristic.}

\ArtifactTable{../artifacts/tables/policy_gap_uncertainty_summary.tex}{Split-level paired gap summary for RC, Austin, and Seattle. The inferential unit is the split-level paired mean net-profit gap; Austin and Seattle remain descriptive because each city has only two split pairs.}

\ArtifactTable{../artifacts/tables/benchmark_bridge_summary.tex}{Benchmark-role bridge summary separating RC's mechanism role from the city impact benchmarks.}

Austin and Seattle preserve real-world-derived geography and travel-time structure but contain only two split pairs per city. They therefore supplement the RC mechanism and uptake-regime studies rather than defining a broad generalization claim. Directionally positive city gaps are useful evidence that the mechanism is not purely an RC artifact, but do not constitute replication or external validation in the inferential sense; stable external validation requires additional city splits and demand calibration.
